
\begin{mdframed}
\textbf{G3 의 위치.} G1(Ch1: 전하보존 $V_n$, $\xi_\eq$ logistic, Level A mobility, spectrum tail,
피팅식 \texttt{eq:fiteq})과 G2(Ch2--4: 발열 연결, ★Level B forward/backward $\chi_j=\beta_j$,
entropy production)를 \emph{수정하지 않고} 적층한다. Ch5 = Level B 의 \textbf{branch 확장}(충방전),
Ch6 = \textbf{수치해법 + 보유 데이터(N-4) 검증}. 지배 표준 GS-1~4 와 10-pass 규율(특히 R3-1
effective-width 정합)을 매 절에 적용.
\end{mdframed}
\newpage

% ====================================================================
\part*{Chapter 5 — 충방전과 히스테리시스 (Level B branch 확장)}
\addcontentsline{toc}{section}{Chapter 5 — 충방전·히스테리시스}
% ====================================================================
\setcounter{section}{0}\renewcommand{\thesection}{5.\arabic{section}}

\begin{anchorbox}
\textbf{관찰과의 끈(N-1).} 관찰의 확장: 같은 SOC 라도 \emph{충전 곡선과 방전 곡선이 어긋난다}.
그러나 그 어긋남(voltage gap)이 모두 ``히스테리시스''는 아니다 — 분극$\cdot$kinetic lag$\cdot$
transport 가 섞인다. Ch5 는 \emph{경로 의존(path dependence)}을 Level B 의 branch 비대칭으로
물리적으로 분리한다.
\end{anchorbox}

\section{목적과 고정 기준}\label{sec:ch5_scope}
Ch1--4 본문(전하보존, $\xi_\eq$, Level A $k_j$, Level B $r_j^\pm$, entropy production)을 유지.
방전 기준 내부 state $\xi_j$ 를 유지하고, 충전 방향 보조 변수는 두되 독립 state 를 늘리지 않는다.
\textbf{원칙}(GS-4): 물리 근거 없는 empirical branch offset$\cdot$arbitrary hysteresis curve 를 본문
기본식으로 쓰지 않는다 — 허용되는 hysteresis 는 metastable phase path, phase-boundary motion,
nucleation barrier, domain memory 로 해석 가능한 경우뿐(\textbf{grounding}: 삽입전극 hysteresis 의
열역학적 기원 = 다입자 metastability/non-monotone single-particle free energy, Dreyer 외
\cite{dreyer2010}; 10-pass R4-G3).

\section{Signed current, branch, signed state coordinate}\label{sec:ch5_signed}
충방전을 한 구조에 넣기 위해 signed current 와 signed 좌표를 둔다:
\begin{equation}
I_s>0:\ \text{방전},\quad I_s<0:\ \text{충전},\quad I_s=0:\ \text{휴지};\qquad
\frac{dq_\state}{dt}=\frac{I_s}{Q_\cell}.
\label{eq:ch5_signed}
\end{equation}
$b(t)\in\{\mathrm{dis,ch,rest}\}$. Ch1--4 의 $|I|$(방전 magnitude)와 $I_s$ 를 구분한다. 전하 보존식은
기준점 포함 상대형(안전):
\begin{equation}
Q_\cell(q_\state-q_{\state,\mathrm{ref}})=\big[Q_\bg(V_n,T)-Q_\bg(V_{n,\mathrm{ref}},T_\mathrm{ref})\big]
+\sum_j Q_{j,\tot}(\xi_j-\xi_{j,\mathrm{ref}}).
\label{eq:ch5_balance}
\end{equation}
충전 방향 표기 $\zeta_j=1-\xi_j$ 는 \emph{종속}(독립 state 아님); 기본 입력은 $\xi_j,\,d\xi_j/dt$.
branch 전환 시 $\xi_j$ 를 재초기화하지 않는다.

\section{True equilibrium vs metastable branch target}\label{sec:ch5_meta}
Ch1 의 $\xi_{\eq,j}(V_n,T)$ 는 \emph{true} thermodynamic equilibrium target. 충방전 경로에서 관측되는
branch 곡선은 metastable phase path$\cdot$nucleation barrier$\cdot$domain memory$\cdot$kinetic lag$\cdot$
분극이 섞인 결과일 수 있다. 물리적으로 해석 가능할 때만 branch free energy $G_j^b$ 로부터 branch
potential $U_j^b(T,h_j)\sim F^{-1}\partial G_j^b/\partial n_j$ 를 정의하고
\begin{equation}
\xi_{j,\tar}^b(V_n,T,h_j)=\Big[1+\exp\!\big(-(V_n-U_j^b)/w_j^b\big)\Big]^{-1}
\label{eq:ch5_branch_target}
\end{equation}
를 \emph{metastable} target 으로 둔다(true eq 아님). 근거 없으면 empirical fit 으로 격하(실무 팁).

\section{★ Branch-dependent Level B kinetics (keystone 확장) 와 정합}\label{sec:ch5_levelB}
Ch3 의 Level B forward/backward 를 branch 별로 둔다:
\begin{equation}
\frac{d\xi_j}{dt}=r_j^{+,b}(1-\xi_j)-r_j^{-,b}\xi_j .
\label{eq:ch5_fb}
\end{equation}
branch free-energy landscape 가 정의되면 그 위에서 \emph{local} detailed balance:
\begin{equation}
\frac{r_j^{+,b}}{r_j^{-,b}}=\frac{\xi_{j,\tar}^b}{1-\xi_{j,\tar}^b}=\exp\!\Big[\frac{V_n-U_j^b}{w_j^b}\Big].
\label{eq:ch5_branch_db}
\end{equation}
이는 \emph{global} equilibrium 의 detailed balance(Ch3 식)가 아니라 metastable branch surface 위의
local balance다.
\begin{keybox}
\textbf{R3-1 정합을 branch 에도 적용(10-pass).} Ch3 R3-1 과 동일하게, branch Level B 의 구동 척도는
branch \emph{thermodynamic} affinity $\mathcal A_{j,\mathrm{prog}}^{b,\mathrm{th}}=RT(V_n-U_j^b)/w_j^b$
이어야 식~\eqref{eq:ch5_branch_db} 와 정합한다(단순 $F(V_n-U_j^b)$ 는 $w_j^b=RT/F$ 극한서만). 즉
forward 장벽 $-\beta_j^b\mathcal A^{b,\mathrm{th}}$, backward $+(1-\beta_j^b)\mathcal A^{b,\mathrm{th}}$.
방전 branch 가 Ch1--3 의 true-eq 극한($U_j^b\to U_j$, $w_j^b\to w_j$, $\xi_\tar^b\to\xi_\eq$)으로
환원되면 G1/G2 와 무충돌(P3-6). \textbf{branch 비대칭의 근원}은 충/방전에서 $U_j^b,w_j^b,\beta_j^b$
(또는 landscape)가 다른 것 — \emph{새 자유 파라미터의 남발이 아니라} metastable branch 의 물리.
\end{keybox}
hysteresis state $h_j\in[-1,1]$ ($+1$ 방전 memory, $-1$ 충전, $0$ relaxed), branch center
$U_j^\hys=U_j^0+\tfrac12\Delta U_j^\hys h_j$, dynamics $dh_j/dt=\lambda_j^b(h_{j,\tar}^b-h_j)$ — $\lambda_j^b$ 는
domain rearrangement$\cdot$phase-boundary depinning$\cdot$nucleation memory time scale(임의 fitting speed
아님; rest relaxation 으로 제약).

\section{관측 voltage gap 의 분해 (히스테리시스 $\ne$ gap)}\label{sec:ch5_gap}
\begin{equation}
\Delta V_\mathrm{obs}=\Delta V_\hys+\Delta V_\pol+\Delta V_\mathrm{kin}+\Delta V_\mathrm{transport}+\Delta V_\mathrm{aging},
\qquad \boxed{\Delta V_\mathrm{obs}\ne\Delta V_\hys.}
\label{eq:ch5_gap}
\end{equation}
branch apparent 전위 $V_{n,\app}^b=V_n+I_s R_n^b$ ($R_n^b>0$ 라도 $I_s$ 부호로 이동 방향 결정 —
polarization 이지 hysteresis offset 아님). full-cell: $V_\cell^b=V_p^b-V_n^b-\eta_\mathrm{loss}^b$, apparent
조합 시 분극 중복 계상 금지.

\section{Hysteresis 에너지와 열}\label{sec:ch5_heat}
loop area $E_\mathrm{loop}=\oint V\,dQ$ 는 true hysteresis + 분극 + kinetic + transport + side
reaction 을 모두 포함 — \textbf{전체를 hysteresis heat 로 단정 금지}. 물리적으로 분리된 hysteresis loss
가 있을 때만(flux--force, $\ge0$):
\begin{equation}
\dot{\mathcal Q}_{\hys,n}^b=J_\hys^b\,\mathcal A_\hys^b\ \ge0\quad(\text{signed }I_s\Delta V_\hys^b\ \text{는 부호 정렬 시만}).
\label{eq:ch5_hysheat}
\end{equation}
branch 가역 열 $\dot{\mathcal Q}_{\rev,\xi}^b=\sum_j s_j^b T\Delta S_j^b (Q_{j,\tot}/F)\,d\xi_j/dt$ (충전서
$d\xi_j/dt<0$ 가능 — 부호 명시); hysteresis heat 와 reversible heat 가 branch 열 비대칭을 \emph{동시에}
설명하지 않도록 독립 제약.

\section{Ch5 검수}\label{sec:ch5_check}
{\small signed current/좌표 우선 정의: 통과. $\xi_j$ state continuity, $\zeta_j$ 종속: 통과.
true eq vs metastable target 구분: 통과. (★R3-1) branch-local detailed balance 가 thermodynamic
affinity $RT(V-U^b)/w_j^b$ 로 정합: 통과. $\Delta V_\mathrm{obs}\ne\Delta V_\hys$: 통과. hysteresis heat
$\ne$ loop area, flux--force $\ge0$: 통과. (P3-6) 방전 true-eq 극한서 G1/G2 환원: 통과. (GS-4)
empirical branch offset 본문 배제: 통과.}

% ====================================================================
\part*{Chapter 6 — 수치해법과 보유 데이터 검증}
\addcontentsline{toc}{section}{Chapter 6 — 수치해법·검증}
% ====================================================================
\setcounter{section}{0}\renewcommand{\thesection}{6.\arabic{section}}

\begin{anchorbox}
\textbf{관찰과의 끈(N-1).} 피팅식(G1 \texttt{eq:fiteq})이 \emph{실제로 풀리고} 보유 데이터로
\emph{검증}되어야 ``저온 긴 꼬리 = 유효배리어 spectrum'' 주장이 닫힌다. Ch6 은 그 solver 구조와
검증 pipeline 을 정한다(P1: 코드 구현은 범위 외, 방법론만).
\end{anchorbox}

\section{목적과 범위}\label{sec:ch6_scope}
Ch1--5 모델을 \emph{음극 방전$\cdot$등온} 1차 검증으로 푼다. 새 물리 추가 없음. 축약/수치법은
(i) 보존법칙 보존, (ii) 극한 거동 보존, (iii) 직접 수치해와 비교검증, (iv) 수학적/문헌 근거 또는
명시적 오차검증 중 하나를 만족해야 한다(GS-4).

\section{전하 보존식 root: dynamic vs OCV}\label{sec:ch6_root}
\textbf{동역학 root}(상태 $\xi_j$ 독립 입력):
$F_\mathrm{cb}^\mathrm{dyn}(V_n;q,\bm\xi)=Q_\bg(V_n)+\sum_j Q_{j,\tot}\xi_j-Q_\cell q=0$,
$\partial F^\mathrm{dyn}/\partial V_n=\partial Q_\bg/\partial V_n$. \textbf{평형 OCV root}($\xi_j=\xi_{\eq,j}(V_n)$):
$\partial F^\OCV/\partial V_n=\partial Q_\bg/\partial V_n+\sum_j Q_{j,\tot}\,\partial\xi_{\eq,j}/\partial V_n$.
\textbf{두 root 의 derivative 가 다르므로 같은 Jacobian 으로 처리 금지}(10-pass R6 부호/도함수 점검 통과).
단조성 $\partial Q_\bg/\partial V_n\ge\epsilon_Q$ 면 해 존재 범위서 root 유일. root 실패는 임의 residual
penalty 로 흡수하지 않고 해당 파라미터 조합을 reject.

\section{Root-constrained ODE / index-1 DAE 시간 적분}\label{sec:ch6_dae}
정전류서 $dq/dt=|I|/Q_\cell$, 상태 $\xi_j$, $V_n$ = algebraic:
\begin{align}
\frac{d\xi_j}{dt}&=k_j^{\mathrm{phys}}(V_n,q,|I|)\,[\xi_{\eq,j}(V_n)-\xi_j]\quad(\text{또는 Level B }r_j^\pm),\\
0&=F_\mathrm{cb}^\mathrm{dyn}(V_n;q,\bm\xi).
\end{align}
순서: $q,\xi_j\to$ root $V_n\to(\xi_\eq,k_j$ 또는 $r_j^\pm)\to d\xi_j/dt\to$ implicit solver(BDF/Radau,
stiff)$\to$ 다음 step; 휴지구간은 $q$ 고정·$\xi_j$ relaxation + root 동시 갱신. \textbf{강구동서 rate 폭주는
임의 cap 으로 자르지 않고 stiff 로 취급}(GS-4), 실제 제한은 물리적 병목($k_{j,\lim}$, transport,
finite current).

\section{★ 분포 적분: Plan B(reference) + Plan A(Refs 6/7 가속, dossier caveat)}\label{sec:ch6_solver}
\textbf{Plan B — direct $g$-grid (항상 보존되는 reference/validator).} barrier 분포 적분
$\xi_j=\int_0^\infty\rho_G(g)\xi_j(g,t)\,dg\approx\sum_m w_m\rho_G(g_m)\xi_j(g_m,t)$. 느릴 수 있으나
가장 직접적인 수치 표현 — \textbf{모든 축약은 이것과 비교 검증}한다.

\textbf{Plan A — 사용자 PhD Refs 6/7 \cite{lee2011,son2013} ratio-substitution(Fredholm/Padé 가속).}
self-consistent kernel integral 을 닫는 \emph{검증 대상 가속 후보}. \textbf{dossier(P3-5) 의
load-bearing 차이를 반드시 반영}:
\begin{boundbox}
\textbf{(R-dossier).} (i) Refs 는 \emph{Fredholm}(고정 구간$\cdot$정상상태); 본 문제는 \emph{Volterra}
(인과, 상한 $q$). Fredholm 환원은 post-peak target 동결 + kernel 국소화 가정이며 그 붕괴는
\emph{측정}한다. (ii) 사용자 논문 자신이 ``reaction zone 이 넓으면 ratio 근사 부정확''이라 경고 —
graphite 에서 \emph{넓은 spectrum = stretched tail(저온$\cdot$관심 영역)} 이므로 \textbf{closure 가 가장
필요한 곳에서 가장 부정확할 수 있다}. $\Rightarrow$ Plan A 단독 채택 금지, Plan B(g-grid) 대비
$\varepsilon$ 를 \emph{운용 $T\cdot$rate$\cdot$tail-window} 에서 측정.
\end{boundbox}
\textbf{승격 조건 F1--F5}: F1 kernel mapping 명확 / F2 $\delta(g-g_0)$ baseline 서 single-barrier
복귀 / F3 Neumann series(또는 ratio truncation) 수렴 수치확인 / F4 0.1--1.0C 서 g-grid 대비
$\varepsilon\le\varepsilon_\mathrm{tol}$ / F5 실패 시 g-grid/adaptive quadrature/Prony 로 fallback.
Prony $K_j(t)\approx\sum_\ell a_\ell e^{-t/\tau_\ell}$ ($a_\ell\ge0,\tau_\ell>0$; 음수 weight/비물리
time constant 금지)은 준정상 근사 — $\xi_j\!\to\!V_n\!\to\!\xi_\eq\!\to\!d\xi_j/dt$ 되먹임 강하면 폐기.

\section{★ 보유 데이터 기반 검증 backbone (N-4)}\label{sec:ch6_validate}
보유: \textbf{음극 반쪽셀 STO--OCV, 반쪽셀 GITT, 0.1/0.2/0.5/1.0C rate series}(등온 가정 유효범위:
자기발열 작을 때). 순차 제한(식별성 보존):
\begin{longtable}{p{0.16\textwidth} p{0.40\textwidth} p{0.38\textwidth}}
\toprule 데이터 & 제한 가능 & 직접 확정 어려움 \\ \midrule \endhead
STO--OCV & 평형 골격 $V_{n,\OCV}(q)$, $U_j,w_j,Q_{j,\tot},Q_\bg$ & $k_j,\rho_G$, transport \\
GITT & instantaneous polarization, relaxation signature, OCV 일관성 & $k_j$ 단독, solid-diffusion vs phase relaxation 완전분리 \\
0.1/0.2C & quasi-eq proxy, peak 위치/폭 & strict equilibrium \\
0.5/1.0C & kinetic lag, broadening, transport stress & pure frozen limit \\
\bottomrule
\end{longtable}
검증 순서: (C1) STO--OCV 로 평형 골격 → (C2) GITT 로 fast polarization·relaxation signature →
(C3) 0.1/0.2C 에서 peak 위치·폭 검증 → (C4) 0.5/1.0C 에서 tail·broadening·shift 검증 →
(C5) 잔차를 $\rho_G,k_{j,\lim},R_n$, transport 중 \emph{독립 제약과 함께} 귀속.

\textbf{★ $\chi_j$ 의 비퇴화 식별(핵심 주장 falsify 설계).} tail 의 $\ln$-Arrhenius slope 가
\emph{전위($V_{n,\drive}-U_j$)에 의존}하는 것이 barrier-spectrum 의 고유 signature(G1
$\partial\ln L/\partial V_\drive=-\chi_j s_{\phi,j}F/RT$). 그러나 charge-transfer 과전압 $\eta_\ct$ 도
Butler--Volmer 라 전위 지수의존 → $\chi_j$ 와 \emph{co-linear} 가능. \textbf{분리 설계}: 다중
$T\times|I|$ + GITT(순간 step 이 $\eta_\ct$, 느린 relaxation 이 barrier-spectrum 을 분리). $\eta_\ct$
독립 분리 \emph{후}에만 $\chi_j$ 귀속 성립. \textbf{Null 규칙}(G1 \S falsify): N1 slope 전위무의존→
transport-dominant 기각; N2 전위/$|I|$ shift 없음→$\chi_j$ 부재; N3 저속 near-peak 저온 미수축→
disorder 분리; N4 $\rho_G$ forward only(역산 금지).

\section{금지되는 피팅 편의와 확장 조건}\label{sec:ch6_forbid}
$\{R_n,k_j^\mathrm{phys},w_j,\rho_G,Q_\bg,Q_{j,\tot},k_{j,\lim}\}$ 를 같은 자료서 \emph{동시 자유화} 금지.
금지: 강구동 rate 를 비물리 cap 으로 절단 / 잔차를 arbitrary background 로 흡수 / hysteresis 를 단순
charge-discharge offset 으로만 / 음수 quadrature weight / Fredholm-Padé correction 을 불분명한 knob 으로.
\textbf{확장 순서}: 음극 방전이 먼저 닫힌 뒤 충전 branch(Ch5 signed/metastable) → 양극(별도 보존식·STO-OCV)
→ full-cell($V_\cell=V_p-V_n-\eta_\mathrm{loss}$, 중복계상 금지).

\section{Ch6 검수}\label{sec:ch6_check}
{\small dynamic vs OCV root 구분(도함수 다름): 통과. 해 존재 조건 residual 흡수 금지: 통과.
DAE/root-constrained ODE 순서: 통과. 강구동 cap 0(stiff+물리 병목): 통과. g-grid = reference,
Plan A = F1--F5 검증 가속(dossier caveat 반영): 통과. GITT 를 $k_j$ 직접값으로 단정 안 함: 통과.
0.1C/1.0C = proxy/stress(strict limit 아님): 통과. $\chi$ vs $\eta_\ct$ 분리 설계 명시: 통과.}

% ====================================================================
\section*{G3 종합 — 6장 합류 완료 + keystone 전파 최종 확인}
\addcontentsline{toc}{section}{G3 종합 — 합류 완료}
% ====================================================================
\begin{longtable}{p{0.08\textwidth} p{0.44\textwidth} p{0.42\textwidth}}
\toprule 장 & 역할(의도 고리) & keystone/검증 정합 \\ \midrule \endhead
Ch1(G1) & 열역학 무대+Level A mobility+spectrum tail+피팅식 & χ=Level A; N-2 \texttt{eq:fiteq} \\
Ch2(G2) & 발열 연결 & Level A→가역=thermo/비가역=flux \\
Ch3(G2) & ★Level B($\chi=\beta$), detailed balance & ★R3-1: $\mathcal A^\mathrm{th}=RT(V-U)/w_j$; $\xi_{\mathrm{ss}}\to\xi_\eq$ 극한 일치 \\
Ch4(G2) & entropy production 발열 & Schnakenberg $(\mathcal J^+-\mathcal J^-)\ln(\cdot)\ge0$ \\
Ch5(G3) & 충방전·hysteresis(Level B branch) & branch-local db 가 $\mathcal A^{b,\mathrm{th}}$ 로 정합; 방전 극한서 G1/G2 환원 \\
Ch6(G3) & 수치해법+보유 데이터 검증 & Plan B reference + Plan A(dossier caveat); $\chi$ vs $\eta_\ct$ 분리 \\
\bottomrule
\end{longtable}
\textbf{합류 완료}: 단일 인과 사슬(N-1) — 관찰 → 열역학 무대($V_n$) → 동역학 주역(Level A mobility,
spectrum tail) → 방향성 확장(Level B, Ch3) → 발열(Ch2,4) → 충방전·hysteresis(Ch5) → 수치·검증(Ch6) —
이 Ch1$\to$Ch6 에 \textbf{keystone 무모순 전파}(P3-6)로 닫혔다. 피팅식(N-2)이 본문 결과로 존재하고,
보유 데이터로 falsify 가능한 검증(N-4)이 설계됐다. \textbf{남은 작업}: G4(xelatex/kotex 빌드 — 사용자
환경) + 사용자 검수(Decision Gate); erf/logistic 최종 판별은 저속 OCV near-peak plot 데이터 대기(DQ-v3-1).

